% Independent proof reconstruction. Not an active book input.
% Compact connected X, Hermitian omega, 2 <= m < n.
% Bounded admissibility uses the chapter's canonical representatives.

\begin{lemma}\label{draft633:bounded-torsion-weight}
Put $k=m-1$ and $\ell=n-m$. For bounded admissible
$t_1,\ldots,t_k$, let
$T=R_k(t_1,\ldots,t_k)\wedge\omega^\ell$.
Then $\ddc T$ is a signed Radon measure. If bounded admissible
$t_i^j\downarrow t_i$ have a common bound, these measures converge
weakly, also after multiplication by any continuous function of
finitely many uniformly bounded decreasing admissible potentials.
\end{lemma}
\begin{proof}
The subset definition and the torsion formulas in
\cref{hessian:bounded-current-torsion} give, in distributions,
\[
\begin{aligned}
 \ddc T={}&\ddc\omega\wedge
 \left(\sum_{a=1}^kR_{k-1}^{\widehat a}\omega^\ell
                       +\ell R_k\omega^{\ell-1}\right)\\
 &+\mathrm d\omega\wedge\dc\omega\wedge
 \left(2\sum_{a<b}R_{k-2}^{\widehat a\widehat b}\omega^\ell
       +2\ell\sum_{a=1}^kR_{k-1}^{\widehat a}\omega^{\ell-1}
       +\ell(\ell-1)R_k\omega^{\ell-2}\right).
\end{aligned}
\]
The sum with two hats is absent when $k=1$, and the last term is
absent when $\ell=1$. Every factor $R_r$ has degree $r\leq k<m$,
so its coefficients are signed measures. Expanding each $R_r$
over subsets expresses the displayed measure as fixed smooth forms
wedged with raw currents $D_p$ of degrees $p\leq m-1$.
\Cref{hessian:rough-coefficient-capacity} therefore bounds its
variation locally by $C c_{m-1}^*$, uniformly under common bounds
on the potentials. Joint convergence follows from
\cref{hessian:joint-decreasing-products} applied to this expansion.
All the canonical weights are quasicontinuous for $c_{m-1}^*$.
\Cref{hessian:weighted-measure-convergence}, with $s=m-1$,
thus proves the weighted assertion coefficientwise on a finite
coordinate cover. A partition of unity gives it on $X$.
The term involving $R_k\ddc(\omega^\ell)$ retains degree $m-1$;
the argument requires no coefficient estimate in degree $m$.
\end{proof}

\begin{proposition}\label{draft633:bounded-nonvanishing}
Every bounded admissible potential $u$ has
\[
 \int_X H_m(u)>0.
\]
No differential identities for powers of $\omega$ are required.
\end{proposition}
\begin{proof}
Suppose instead that $H_m(u)=0$. Fix $s>0$, set
$c=\inf_Xu+s$, and put $w=\max\{u,c\}$.
The upper semicontinuity of $u$ makes
$U=\{u<c\}$ open, and the definition of the infimum makes it
nonempty. Moreover $0\leq w-u\leq s$ and this difference is
zero outside $U$. The maximum is bounded admissible with its
canonical pointwise values by \cref{hessian:regularized-maximum}.

For $0\leq q\leq m$, define the positive Radon measures
\[
 \mu_q=Q_m(\underbrace{u,\ldots,u}_{q},
                    \underbrace{w,\ldots,w}_{m-q}).
\]
Thus $\mu_m=H_m(u)=0$. For $q<m$, let $T_q$ be the sum of
$m-q$ currents with $q+r$ slots equal to $u$ and
$m-q-1-r$ slots equal to $w$, for $0\leq r<m-q$, each of degree
$m-1$ and wedged with $\omega^{n-m}$. In smooth notation this is
\[
 T_q=\sum_{r=0}^{m-q-1}
       (\omega+\ddc u)^{q+r}
       (\omega+\ddc w)^{m-q-1-r}\omega^{n-m}.
\]
For bounded potentials this notation means the symmetric shifted
current defined by the subset expansion, not an unproved ordinary
product of distributional factors.

Choose smooth bounded admissible $u_j\downarrow u$ and the
regularized maxima $w_j=M_{\epsilon_j}(u_j,c)\downarrow w$,
where $\epsilon_j\downarrow0$, as in
\cref{hessian:regularized-maximum}.
Both lists have a common bound. The smooth polynomial identity and
Stokes' theorem on the compact manifold give
\[
 \int_X(\mu_q^j-\mu_m^j)
              =\int_X(w_j-u_j)\ddc T_q^j.
\]
Joint decreasing convergence passes the left side to its limit.
The preceding lemma, with $F(x,y)=y-x$ and the decreasing lists
$u_j,w_j$, passes the right side to its limit. Hence
\[
 \int_X\mu_q=\int_X(w-u)\ddc T_q.
\]
The difference of the two decreasing sequences need not itself
decrease; the continuous two-variable weight is essential here.

On the open set $U$, the potential $w$ is constant. Locality of
the raw recursion therefore sets every factor $\ddc w$ to zero
there, so that its shifted factors become $\omega$.
The formula in the preceding lemma expresses $\ddc T_q|_U$ as
finitely many smooth metric-derivative forms wedged with powers
of $\omega+\ddc u$ of degrees at most $m-1$.
Use smooth admissible approximations to $u$ and the constant $c$ in
these local expressions; their limits are the same restricted currents.
For those smooth approximations, the absolute mixed-form
estimate \cref{hessian:absolute-mixed-forms} bounds each such term
by a constant times the corresponding positive trace, and therefore
by $C\sum_{r=0}^{m-1}\mu_r$ on $U$ after taking joint limits.
More explicitly, a term with $r$ factors is bounded by
$C(\omega+\ddc u)^r\omega^{n-r}$; this is $C\mu_r$ on $U$.
Against compactly supported nonnegative tests in $U$, weak
convergence and lower semicontinuity of total variation preserve
the estimate. Exhaustion gives the measure inequality
\[
 \mathbf1_U|\ddc T_q|
                 \leq C\sum_{r=0}^{m-1}\mathbf1_U\mu_r.
\]
Here $C$ depends only on $\omega,m,n$, not on $u$ or $s$; a fixed
finite coordinate cover bounds all the smooth metric coefficients.

Since $w-u$ vanishes outside $U$, positivity gives
\[
 \int_U\mu_q\leq\int_X\mu_q
            \leq Cs\sum_{r=0}^{m-1}\int_U\mu_r
                 \qquad(0\leq q<m).
\]
Summing these $m$ inequalities and choosing $s>0$ with $mCs<1$
forces $\sum_{r=0}^{m-1}\int_U\mu_r=0$.
But locality gives $\mu_0=\omega^n$ on the nonempty open set $U$,
whose volume is positive. This contradiction proves the result.
\end{proof}
